\chapter{Геометричні ефекти}\label{ch:geometry}

Рівновага (розд.~\ref{ch:equilibrium}) дає функцію $\psi(R,Z)$. У цьому розділі з неї будуємо координати, в яких записують рівняння переносу й потоків (розд.~\ref{ch:transport}, \ref{ch:flows}), розбираємо, як форма перерізу впливає на стійкість, і закінчуємо топологією $\psi$: критичними точками, сепаратрисою та дослідницьким треком проєкту.

\section{Системи координат і метрика}\label{sec:04-coords}

\begin{outofcorpus}[потокові координати, кут прямих силових ліній, усереднення]
Візьмемо $(\psi,\theta,\tor)$, де $\theta$ --- будь-який полоїдальний кут на поверхні $\psi=\text{const}$, і якобіан $\mathcal{J}^{-1}=\grad\psi\times\grad\theta\cdot\grad\tor$ (об'єм $\dd V=\mathcal{J}\,\dd\psi\,\dd\theta\,\dd\tor$). З $\Bvec=\grad\psi\times\grad\tor+F\grad\tor$ \eqref{eq:03-B} маємо $\Bvec\cdot\grad\tor=F/R^2$ і $\Bvec\cdot\grad\theta=-1/\mathcal{J}$. Тому лінія поля має $\dd\theta/\dd\tor=-R^2/(F\mathcal{J})$, а запас стійкості та кут прямих силових ліній (straight-field-line angle)
\begin{equation}\label{eq:04-q}
q(\psi)=\frac{1}{2\pi}\oint\frac{F|\mathcal{J}|}{R^2}\,\dd\theta=\frac{\dd\psit}{\dd\psi},\qquad
\thetastar=\frac{1}{q}\int_0^{\theta}\frac{F|\mathcal{J}|}{R^2}\,\dd\theta',\qquad \frac{\dd\thetastar}{\dd\tor}=\frac1q
\end{equation}
(знак залежить від орієнтації кутів і напрямків $\Ip$, $\Btor$). Друга рівність для $q$ випливає з того, що тороїдальний потік на радіан $\psit=\Psit/2\pi=\frac{1}{2\pi}\int\Btor\,\dd A$ ($\Psit$~--- повний потік), а елемент площі перерізу $\dd A=|\mathcal{J}|R^{-1}\dd\psi\,\dd\theta$. У змінних $(\thetastar,\tor)$ силові лінії --- прямі; саме ця пара разом з $\psit$ утворює канонічні змінні з розд.~\ref{ch:particles}.

\emph{Усереднення по магнітній поверхні}: $\fsa{f}=\oint f\,|\mathcal{J}|\,\dd\theta\big/\oint|\mathcal{J}|\,\dd\theta$, $V'=\dd V/\dd\psi=2\pi\oint|\mathcal{J}|\,\dd\theta$. Основні тотожності для однозначних осесиметричних $f$ (або з усередненням і по $\tor$) і векторних полів $\vect{A}$:
\begin{equation}\label{eq:04-fsa}
\fsa{\Bvec\cdot\grad f}=0,\qquad \fsa{\grad\cdot\vect{A}}=\frac{1}{V'}\frac{\dd}{\dd\psi}\Big(V'\fsa{\vect{A}\cdot\grad\psi}\Big).
\end{equation}
Перша випливає з $\Bvec\cdot\grad f=-\mathcal{J}^{-1}\pa_\theta f+(F/R^2)\pa_\tor f$ (інтеграл від повної похідної по замкненому контуру дорівнює нулю). Друга --- з формули дивергенції в координатах. Саме \eqref{eq:04-fsa} прибирає з рівнянь переносу паралельну динаміку й залишає одновимірні рівняння за $\psi$ (розд.~\ref{ch:transport}).
\end{outofcorpus}

\paragraph{Координати Бузера.} Це координати прямих силових ліній, у яких ще й коваріантні компоненти $B_\theta$, $B_\tor$ є функціями потоку (стандартне означення, поза корпусом). Тищенко~\cite[с.~25]{Tyshchenko2021} записує рівняння альфвенівських хвиль у координатах Бузера $(\Psit,\theta,\tor)$ з радіальною змінною~--- тороїдальним потоком і поздовжнім оператором $q^{-1}\pa_\theta+\pa_\tor$ (у джерелі полоїдальний кут позначено іншою літерою, а в підписі до (1.4) назви кутів переставлено; з множника $\exp(im\theta-in\tor)$ видно, що $\theta$ --- полоїдальний). Метрика в параксіальному наближенні~\cite[с.~26]{Tyshchenko2021}: $B=\bar B[1+\epsilon_B\cos(m_c\theta-n_cN\tor)]$ і аналогічно для $g^{ij}$, причому $g^{11},g^{22}\propto\delta_\kappa=(\kappa+\kappa^{-1})/2$. «Параметри зачеплення» $\epsilon_B$, $\epsilon_g$ зв'язують гармоніки $m$ і $m+m_c$: $m_c=1$ відповідає тороїдальності, $m_c=2$ --- витягнутості; у токамаку $n_c=0$. Отже, форма перерізу прямо входить у спектр хвиль через метрику.

\paragraph{MSCS Seto.} Для 2D-моделі ядра й периферії Seto і Fukuyama~\cite[с.~2]{Seto2014} беруть систему координат магнітної поверхні (MSCS). У наших позначеннях це $(\rho_{\mathrm{M}},\theta_{\mathrm{M}},\tor)$ (в оригіналі кути позначено іншими літерами). Радіальна мітка $\rho_{\mathrm{M}}$ будується з тороїдального потоку; полоїдальний кут $\theta_{\mathrm{M}}$ --- нормована довжина проєкції силової лінії на площину $\tor=\text{const}$; $\tor$ --- геометричний тороїдальний кут, тому $\vect{e}_\tor=R^2\grad\tor$ ортогональний до двох інших напрямків. Поле $\Bvec=\grad\tor\times\grad\psi+I\grad\tor$ (наше $\psi$ з точністю до знака і нормування, $I\equiv F$), $q=\dd\psit/\dd\psi$ у згоді з \eqref{eq:04-q}. Оскільки кути визначено незалежно від потокових функцій, MSCS «не потокова» і, на відміну від $(\psi,\thetastar,\tor)$, придатна й за сепаратрисою, де $q\to\infty$.

\paragraph{Координати GTC.} Гірокінетичний код GTC~\cite[с.~6--7]{Wang2006} використовує потокові координати прямих силових ліній з геометричним $\tor$ і радіусом $r=\sqrt{\Psit/\Psi_{\mathrm{t,b}}}$ ($\Psi_{\mathrm{t,b}}$ --- на межі плазми). Масштаб дрейфової турбулентності $\propto\rho_{\mathrm i}\propto\sqrt{T_{\mathrm i}}$, а в NSTX $T_{\mathrm i}$ падає від $\sim$\,кеВ у центрі до $\sim\SI{10}{еВ}$ біля сепаратриси~\cite[с.~6]{Wang2006}. Тому сітку роблять рівномірною за
\begin{equation}\label{eq:04-gtc}
\frac{\dd\rho_{\mathrm{G}}}{\dd r}=\sqrt{\frac{T_{\mathrm c}}{T_{\mathrm i}(r)}},\qquad \Delta r\sim\sqrt{T_{\mathrm i}(r)/T_{\mathrm c}},
\end{equation}
де $T_{\mathrm c}$ --- температура в опорній точці~\cite[с.~7, (1)]{Wang2006}. Полоїдальний крок $\Delta\theta(r)$ сталий на поверхні і добирається так, щоб дуга біля екватора була $\sim\rho_{\mathrm i}$; тривимірна сітка йде вздовж наближених ліній $\bar q(r)\theta-\tor=\text{const}$.

\section{Форма перерізу і вертикальна стійкість}\label{sec:04-shape}

\begin{outofcorpus}[параметри форми]
За крайніми точками LCFS: $R_{\mathrm{geo}}=(R_{\max}+R_{\min})/2$, $a=(R_{\max}-R_{\min})/2$, аспектне відношення $A=R_{\mathrm{geo}}/a$, витягнутість $\kappa=(Z_{\max}-Z_{\min})/2a$, верхня і нижня трикутності $\dup=(R_{\mathrm{geo}}-R_{Z_{\max}})/a$, $\dlo=(R_{\mathrm{geo}}-R_{Z_{\min}})/a$, де $R_{Z_{\max}}$ --- радіус найвищої точки. Витягнутість дає більший струм за того самого $q$, але потребує зовнішнього поля з від'ємним індексом спаду $n_{\mathrm d}<0$, що робить плазму вертикально нестійкою.
\end{outofcorpus}

\paragraph{Індекс спаду поля.} Для зовнішнього вертикального поля $n_{\mathrm d}=-(R/\BZ)\,\pa\BZ/\pa R$. Кругова плазма вертикально стійка при $n_{\mathrm d}>0$ і горизонтально стійка при $n_{\mathrm d}<3/2$~\cite[с.~76]{Merlo2011}. У double-null (DN) конфігурації RFX-mod котушки F8 збільшують об'єм і трикутність, витягуючи плазму назовні, і на різних ділянках екватора порушено \emph{обидві} умови. Лінеаризована модель з резистивними пасивними структурами (ще без сідлових котушок) має дві нестійкі моди, причому горизонтальна домінує: $\tau_{\mathrm{hor}}=\SI{51.12}{мс}$, $\tau_{\mathrm{vert}}=\SI{60.72}{мс}$~\cite[с.~75--76]{Merlo2011}. Активну стабілізацію розглянуто в розд.~\ref{ch:stability}.

\begin{corpusexample}[пасивна стабілізація і деформованість~\cite{Ward1992}]
NOVA-W розраховує вертикальну моду з двома симетричними провідними пластинами (кожна $\approx1/8$ периметра стінки, опір як в алюмінію товщиною \SI{1}{см}); кут $|\Theta_{\mathrm{pl}}|$ відраховують від зовнішнього екватора (с.~11).
\begin{itemize}
\item Еліпс $A=100$, $\kappa=2$: оптимум $|\Theta_{\mathrm{pl}}|=\pi/2$ (над і під плазмою), інкремент лише $\approx3{,}5$ раза більший, ніж з повною стінкою (с.~11--12).
\item Еліпс $A=4{,}5$ з витягнутістю ARIES-I і нульовою трикутністю: оптимум зсувається назовні, $|\Theta_{\mathrm{pl}}|\approx0{,}448\pi$ (тороїдальність, с.~15).
\item CIT ($R_0=\SI{2.182}{м}$, $a=\SI{0.660}{м}$, $\kappa_{95}=1{,}996$, $\delta_{95}=0{,}258$, с.~17): $|\Theta_{\mathrm{pl}}|\approx0{,}403\pi$; ідеальна стійкість лише при $\pi/4\lesssim|\Theta_{\mathrm{pl}}|\le\pi/2$ (с.~18--19).
\item ARIES-I ($\kappa=1{,}61$, $\delta=0{,}43$): $|\Theta_{\mathrm{pl}}|\approx0{,}394\pi$, хоча $A$ більше, ніж у CIT. Трикутність зсуває оптимум назовні сильніше, ніж тороїдальність (с.~19).
\end{itemize}
Причина --- деформованість: плазма «обтікає» пластини, гармоніки $m=2,3$ перебудовуються і зменшують стабілізуючі вихрові струми. Рівноваги в~\cite{Ward1992} штучно дестабілізовано, тож числа ілюструють тенденцію, а не проєкт.
\end{corpusexample}

Профіль струму впливає на форму через $\li$ і $\Lambda$ \eqref{eq:03-Bv}. За стандартною теорією (поза корпусом), що більша $\li$, то менше струму на периферії і то важче стабілізувати витягнутий переріз (розд.~\ref{ch:stability}).

\section{Конфігурації}\label{sec:04-configs}

\textbf{Лімітер і дивертор.} LCFS або торкається лімітера, або проходить крізь X-точку, і тоді межа --- сепаратриса~\cite[с.~30]{Merlo2011}. Дивертор відводить найбільші теплові навантаження в точки перетину сепаратриси зі стінкою, керовані положенням X-точки (розд.~\ref{ch:wall}). \textbf{Double-null} (дві X-точки): у RFX-mod розряд \#29746 відтворено з $\li=0{,}95$, $\betap=0{,}29$~\cite[с.~38]{Merlo2011}; реконструкція межі підтверджує, що DN було досягнуто, але згодом конфігурація деградувала до лімітерної, бо регулятора форми не було~\cite[с.~54]{Merlo2011}. \textbf{Від'ємна трикутність}: AUG \#36026 --- перша спроба на AUG, переважно через верхню трикутність. Fenix добре відтворює розряд до \SI{3}{с}, але в H-моді симуляція дає плазму, обмежену \emph{зовнішнім} лімітером, а реконструкція --- \emph{внутрішнім}; причина невідома~\cite[с.~7--8]{Fable2022}. \textbf{Сферичний токамак} NSTX: $A\approx1{,}2$ (модельні $R_0$, $a$ з табл.~\ref{tab:01-devices}). TSC+DCON прогнозують, що при нагріванні NBI до \SI{5}{МВт} енергетично досяжне $\beta>15\,\%$, але балонна нестійкість починається при $\beta\approx12\,\%$ через малий шир біля $q=1$ і пікований тиск~\cite[с.~5]{Jardin2000}. \textbf{Симетрія і власне обертання.} Напрям спонтанного обертання визначає симетрія поля (токамак проти геліотрона)~\cite[с.~6--7]{Ida2001}; див. п.~\ref{sec:06-intrinsic}.

\textbf{Гофрування TF-поля} --- порушення осесиметрії з періодом $2\pi/N_{\mathrm{TF}}$ --- розглянуто в п.~\ref{sec:02-ripple}. Усе нижче стосується осесиметричного $\psi$.

\section{Топологія $\psi(R,Z)$: O- і X-точки, сепаратриса}\label{sec:04-topology}

\begin{figure}[htbp]\centering
\begin{minipage}[c]{0.38\textwidth}\centering
\begin{tikzpicture}[scale=0.62,>=Stealth]
 \foreach \s in {0.35,0.6,0.85}
   \draw[gray] (0.08,0.15) ellipse ({1.75*\s} and {2.5*\s});
 \draw[thick,blue!70!black] (0,-2.6) .. controls (1.3,-1.7) and (2.0,-0.6) .. (2.0,0.4)
   .. controls (2.0,2.2) and (1.0,2.9) .. (0,2.9) .. controls (-1.0,2.9) and (-2.0,2.2) .. (-2.0,0.4)
   .. controls (-2.0,-0.6) and (-1.3,-1.7) .. (0,-2.6);
 \draw[thick,blue!70!black] (0,-2.6) -- (0.75,-3.4) (0,-2.6) -- (-0.75,-3.4);
 \draw[dashed,gray] (-1.3,-3.4) .. controls (-2.6,-1.5) and (-2.6,3.4) .. (0,3.45)
   .. controls (2.6,3.4) and (2.6,-1.5) .. (1.3,-3.4);
 \draw[very thick] (-1.6,-3.4) -- (1.6,-3.4);
 \fill (0.08,0.15) circle (2.5pt) node[right,font=\scriptsize] {O ($+1$)};
 \fill[red!70!black] (0,-2.6) circle (2.5pt) node[right=2pt,font=\scriptsize] {X ($-1$)};
 \node[font=\scriptsize,blue!70!black] at (2.05,2.6) {LCFS};
 \node[font=\scriptsize,gray] at (-2.35,2.95) {SOL};
\end{tikzpicture}
\end{minipage}\hfill
\begin{minipage}[c]{0.6\textwidth}\centering
\begin{tikzpicture}\node[inner sep=0] (im) {\includegraphics[width=\textwidth,trim=405 168 405 205,clip]{viz_plasma.jpg}};
 \fill[black] ($(im.north east)+(-7mm,0)$) rectangle ($(im.north east)+(0,-11mm)$);\end{tikzpicture}
\end{minipage}
\caption{Ліворуч: схема критичних точок $\psi$ в однонульовій конфігурації (у дужках --- індекс Пуанкаре--Хопфа), сепаратриса (LCFS) з двома ногами до дивертора, відкриті лінії SOL. Праворуч: об'ємний рендер плазми DIII-D \#203702, $t=\SI{1760}{мс}$, $\qnf=3{,}70$ (фрагмент рис. \texttt{viz\_plasma}); колір --- $\psiN$ від осі до краю, блакитні лінії --- силові лінії на магнітних поверхнях. Власний розрахунок~\cite{ProjViz}; дані рівноваги: Fusion Equilibrium Challenge (Sophelio / General Atomics), CC BY 4.0.}\label{fig:04-topo}
\end{figure}

\begin{outofcorpus}[критичні точки, індекс, теорія Морса]
Критична точка $\grad\psi=0$ --- це нуль $\Bpol=|\grad\psi|/R$. Гесіан $\mathsf{H}_\psi$ класифікує її: $\det\mathsf{H}_\psi>0$ --- еліптична O-точка (магнітна вісь або центр острова), індекс $+1$; $\det\mathsf{H}_\psi<0$ --- гіперболічна X-точка, індекс $-1$. Індекс дорівнює числу обертів вектора $\grad\psi$ уздовж малого контуру навколо точки. \emph{Теорема Пуанкаре--Хопфа}: якщо $\grad\psi$ трансверсальний до межі області $D$, сума індексів усередині дорівнює ейлеровій характеристиці $\chi_{\mathrm{E}}(D)$; для диска $\chi_{\mathrm{E}}=1$. \emph{Теорія Морса}: топологія підрівневої множини $\{\psiN\le t\}$ змінюється лише при проходженні критичних значень; мінімум додає компоненту ($\chi_{\mathrm{E}}\to\chi_{\mathrm{E}}+1$), сідло --- з'єднує дві компоненти або додає дірку ($\chi_{\mathrm{E}}\to\chi_{\mathrm{E}}-1$). Для вкладених поверхонь (одна O-точка, всередині LCFS сідел немає) $\chi_{\mathrm{E}}(\{\psiN\le t\})=1$ при всіх $t<1$. Сепаратриса --- рівень $\psi$, що проходить через X-точку; у диверторній плазмі вона і є LCFS. Шум, який народжує пару O--X, не змінює суми індексів.
\end{outofcorpus}

Критичні точки --- це саме те, що рахують скаляри рівноваги: положення осі, форма межі через сідло (розд.~\ref{ch:equilibrium}, E1--E3). Проєкт перевірив, чи дає скелет критичних точок (рис.~\ref{fig:04-topo}) незалежний критерій якості передбаченого $\psi$~\cite{ProjRegisters,FEC2026}. Індекс рахують як число обертів $\grad\psi$ навколо 8 сусідів кожного пікселя. Сусідні пікселі з ненульовим індексом об'єднують у компоненту, і для неї рахують число обертів по периметру обмежувального прямокутника (одна справжня точка займає блок $2\times2$). Область --- внутрішність LCFS, ерозована на 2 пікселі, щоб X-точки не потрапили в неї.

\begin{established}{E9--E11}
\cite{ProjRegisters}; \texttt{topology\_probe.py}, DIII-D \#203702, 8 кадрів, гаусів шум $\varepsilon\cdot\mathrm{std}(\psi)$ (перевірено повторним прогоном):
\begin{center}\small
\begin{tabular}{@{}lcccccc@{}}\toprule
$\varepsilon$ & 0 & 0,001 & 0,003 & 0,01 & 0,03 & 0,10\\\midrule
$R^2_\psi$ & 1 & 1 & 0,99999 & 0,99993 & 0,99935 & 0,99278\\
O / X (медіана) & 1 / 0 & 1 / 0 & 1 / 0 & 1 / 0 & 1 / 0 & 10,5 / 11,0\\
сума індексів & $+1$ & $+1$ & $+1$ & $+1$ & $+1$ & $+1{,}5$\\\bottomrule
\end{tabular}
\end{center}
E9: істина має рівно одну O-точку і жодного сідла. E10: скелет стійкий до 3\,\% шуму, а на 10\,\% з'являється 20,5 зайвих точок. E11: сума індексів при цьому лишається близькою до $+1$, бо зайві точки народжуються парами O--X. Тому інформативна лише \emph{кількість} точок, а не сума. Кожна зайва O-точка --- острів, якого немає.
\end{established}

\begin{established}{E12--E13}
\texttt{topology\_ec\_curve.py}: на істині $\chi_{\mathrm{E}}(\{\psi\le t\})=1$ на всіх 40 рівнях усередині LCFS (кожна підрівнева множина --- диск; еталон без налаштування). Крива $\chi_{\mathrm{E}}(t)$ менш чутлива, ніж лічильник: при 3\,\% шуму відхилення $\overline{|\Delta\chi_{\mathrm{E}}|}=0{,}14$, а медіана кількості точок 1,5 («лічильник спрацював»; у \texttt{topology\_ec\_curve.py} інша реалізація шуму; на 3\,\% лічильник на межі спрацювання: 1 або 1,5); при 10\,\% --- 4,1 і 17. За Морсом крива є інтегралом від лічильника~\cite{ProjRegisters}.
\end{established}

\begin{refuted}{R5}
«Крива $\chi_{\mathrm{E}}(t)$ краща за лічильник критичних точок». Ні: вона менш чутлива (E13), бо за Морсом лише підсумовує ту саму інформацію; пара мінімум--сідло, чиї критичні значення лежать між двома сусідніми рівнями $t$, не змінює $\chi_{\mathrm{E}}$ на жодному рівні (внески $+1$ і $-1$ скорочуються)~\cite{ProjRegisters}. Роль кривої --- бінарний шлюз допустимості, а не градуйована метрика.
\end{refuted}

\begin{established}{E15--E16}
\texttt{h5\_transfer\_topology.py}: DIII-D і MAST обидві мають скелет 1O/0X із сумою $+1$ (8/8 кадрів кожна), попри протилежні конвенції знака (\texttt{AXIS\_SIGN}: DIII-D $-1$, MAST $+1$), різні сітки й геометрію --- критерій машинно-незалежний (E15). Проєкція MAST на PCA-базис DIII-D з $k=5$, 10, 20, 50 компонентами дає $R^2_\psi=0{,}66$, $0{,}88$, $0{,}93$, $0{,}96$, але скелет лишається 1O/0X (E16)~\cite{ProjRegisters}. Уточнення 2026-09-29: при $k=50$ позначка «1O/0X» у таблиці скрипта --- це медіана $N_X=0{,}5$, округлена до цілого; у 4 із 8 кадрів є внутрішня X-точка. Твердження E16 для $k=5$ це не зачіпає.
\end{established}

\begin{refuted}{R7}
Гіпотеза H5: «провал перенесення DIII-D$\to$MAST має топологічну складову». Спростовано E16: навіть при $R^2_\psi=0{,}66$ топологія ціла. Провал кількісний, а не топологічний~\cite{ProjRegisters}.
\end{refuted}

\begin{refuted}{R9 (була гіпотеза C8)}
\textbf{Було:} збої вилучення LCFS при перенесенні (3/8 кадрів при $k=5$ і 2/8 при $k=20$) --- систематичний сигнал. \textbf{Спростовано 2026-09-29:} усі збої припадають на кадри, де підбір знака обрав $-1$. Функція \texttt{skeleton()} викликала \texttt{extract\_lcfs} без \texttt{axis\_sign}, і оцінювач переходив до резервного варіанта «перший успішний». Якщо повернути реконструкції знак MAST, збоїв немає за жодного $k$. Отже, це артефакт знака~$\psi$, а не властивість перенесення. Скрипт перевірки: \texttt{Our try/04-novelty/c8\_sign\_check.py}~\cite{ProjRegisters}.
\end{refuted}

\begin{practicum}[скелет критичних точок]
Дані: 6 демо-розрядів (3 DIII-D, 3 MAST, 67~МБ) уже лежать у \texttt{fusion equilibrium challenge/starter/parquet\_data/}; окремо їх можна завантажити з GitHub starter kit (Git LFS). Середовище --- \texttt{starter/.venv} (системний \texttt{python3} без \texttt{numpy} не підійде). З теки \texttt{Our try/04-novelty}:\\
\texttt{PY="../../fusion equilibrium challenge/starter/.venv/bin/python"}\\
\texttt{"\$PY" topology\_probe.py; "\$PY" topology\_ec\_curve.py; "\$PY" h5\_transfer\_topology.py}\\
Кожен скрипт працює $\approx1$~с і має відтворити таблицю E9--E11 та числа E12--E16. \textbf{Завдання:} (1)~для 10 реалізацій шуму знайдіть частку, де медіана кількості точок $>1$, при 2, 3, 5\,\%; (2)~вимкніть ерозію (\texttt{erode=0}) і перевірте засторогу проєкту: якщо X-точки потрапляють в область, сума стає $1-k_X$, і критерій ламається на double-null; (3)~замініть білий шум гладким (гаусове згладження) --- чи бачить його скелет?
\end{practicum}

\subsection*{Контрольні запитання}
\begin{enumerate}
\item Виведіть $q=\dd\psit/\dd\psi$ з \eqref{eq:04-q}. Чому в координатах $(\thetastar,\tor)$ силові лінії прямі?
\item Чому MSCS~\cite{Seto2014} придатна за сепаратрисою, а $(\psi,\thetastar,\tor)$ --- ні?
\item Навіщо GTC нерівномірна радіальна сітка \eqref{eq:04-gtc}? Що станеться з рівномірною сіткою в NSTX?
\item Чому оптимальні пасивні пластини зсуваються назовні зі зростанням тороїдальності й трикутності~\cite{Ward1992}?
\item Чому сума індексів не реагує на шум, а кількість критичних точок реагує? Чим відрізняються ролі лічильника і кривої $\chi_{\mathrm{E}}(t)$?
\end{enumerate}

\subsection*{Задачі}
\begin{enumerate}
\item \emph{Пластини CIT.} Для CIT (п.~\ref{sec:04-shape}) обчисліть $A$ і, в еліптичному наближенні $R=R_0+a\cos\Theta_{\mathrm{pl}}$, $Z=\kappa a\sin\Theta_{\mathrm{pl}}$, координати центру оптимальної пластини при $|\Theta_{\mathrm{pl}}|=0{,}403\pi$ і при $0{,}448\pi$. На скільки градусів зсунуто оптимум від $\pi/2$?
\item \emph{Сітка GTC.} Нехай $T_{\mathrm i}(r)=T_{\mathrm c}(1-r^2)$. Знайдіть $\rho_{\mathrm G}(r)$ з \eqref{eq:04-gtc} і частку вузлів рівномірної сітки за $\rho_{\mathrm G}$, що припадає на $r>0{,}9$. У скільки разів радіальний крок біля сепаратриси NSTX ($T_{\mathrm i}\approx\SI{10}{еВ}$) менший, ніж у центрі ($\approx\SI{1}{кеВ}$)?
\item \emph{Пуанкаре--Хопф.} (а)~У диску всередині LCFS шум створив 20 зайвих критичних точок. Скільки серед них O і скільки X? (б)~Область-диск містить рівно одну O- і одну X-точку. Чому тоді $\grad\psi$ не може бути трансверсальним до межі, і що це означає для вибору області в \texttt{topology\_probe.py}?
\end{enumerate}
